\section{Conclusion}

Transit accessibility edge weights provide a verifiable, non-partisan mechanism for keeping transit communities together in algorithmic redistricting. The signal respects the urban-rural divide by treating both transit-rich urban tracts and transit-poor rural tracts as internally similar, while identifying the transition boundary (urban transit zone edge) as a natural cut point.

The 1--1.5\% compactness cost is small relative to the community-of-interest benefit, and the signal degrades gracefully in states with limited GTFS coverage. Transit weights are particularly valuable in large metro areas (Charlotte, Milwaukee, Houston, Austin, Dallas) where transit corridors define community identity and political interests in ways not captured by topographic, economic, or administrative signals alone.

Combined with the other six M-track signals, transit accessibility contributes to the composite Community Character Index (M.8) as the infrastructure-access component of community identity.
